\chapter{Justificación}

La creación de una revista científica interna para la carrera de Ingeniería en Biotecnología se fundamenta en la creciente evidencia sobre los beneficios de las publicaciones estudiantiles en la formación profesional. Diversos estudios han documentado que las revistas dirigidas por estudiantes constituyen un espacio formativo de alto valor pedagógico, al tiempo que contribuyen al ecosistema de investigación académica \citep{warren2025vital}.

En primer lugar, la experiencia de publicación permite a los estudiantes familiarizarse con el proceso editorial y de revisión por pares, competencias que tradicionalmente se adquieren hasta etapas avanzadas del posgrado. Como señalan Antony y colaboradores \citep{antony2025facilitating}, la participación en revistas estudiantiles facilita el aprendizaje del proceso de revisión por pares y promueve el desarrollo del pensamiento crítico, la alfabetización informacional y las habilidades de colaboración. De manera similar, la experiencia de la Nazarbayev University Graduate School of Education muestra que los estudiantes que participan como autores o revisores en revistas internas adquieren confianza para someter sus manuscritos a publicaciones internacionales y desarrollan una visión de la autoría como participación en una comunidad de práctica \citep{drybrough2025eight}.

En segundo lugar, la autenticidad de la audiencia constituye un factor determinante en la motivación y el desempeño de los estudiantes en tareas de escritura académica. Investigaciones en el ámbito de la formación docente han demostrado que escribir para una audiencia real—y no únicamente para la evaluación del docente—mejora significativamente la motivación, la autoeficacia y la percepción general de la escritura académica \citep{banegas2020learning}. Este hallazgo es consistente con lo reportado por la literatura sobre publicaciones estudiantiles, que indica que los estudiantes toman los proyectos más seriamente y aprenden más cuando el trabajo está dirigido a una audiencia externa en lugar de a una tarea de clase \citep{uaa2014undergraduate}. En el contexto específico de la escritura para publicación, Saputri y colaboradores \citep{saputri2023students} encontraron que la motivación inicial de los estudiantes noveles por publicar es alta, aunque puede fluctuar ante revisiones extensas, lo que subraya la importancia de contar con acompañamiento y procesos editoriales claros.

En tercer lugar, las revistas estudiantiles contribuyen a la construcción de la identidad académica de los participantes. Un estudio autoetnográfico sobre editores de revistas estudiantiles encontró que la experiencia editorial es relevante para el desarrollo de la identidad académica, particularmente en aspectos como el pensamiento reflexivo, la identidad autoral, la confianza y el aprendizaje a través de la crítica \citep{lujano2024exploring}. Este proceso es especialmente valioso en etapas tempranas de la formación, pues permite a los estudiantes visualizarse como productores de conocimiento y no solo como consumidores del mismo.

Las revistas científicas dirigidas por estudiantes han ganado reconocimiento como un elemento vital del ecosistema de investigación, con experiencias documentadas en diversas disciplinas como ciencias, medicina, derecho, políticas públicas y filosofía, entre otras \citep{warren2025vital}. El Student Journal Symposium, que en su segunda edición reunió a 73 estudiantes y profesionales representando 34 revistas estudiantiles de 28 universidades, evidencia el creciente interés y la consolidación de este tipo de iniciativas a nivel internacional \citep{warren2025vital}. Asimismo, la experiencia de la revista \textit{Current Issues in Education}, que celebra 25 años como publicación dirigida por estudiantes de posgrado, demuestra que estos proyectos pueden alcanzar sostenibilidad y generar impacto académico más allá de su institución de origen \citep{lujano2023special}.

Por otra parte, Huard \citep{huard2019showcase} distingue entre revistas de exhibición, que publican sin revisión por pares, y revistas selectivas, que implementan un proceso formal de revisión. La autora concluye que, si bien cualquier plataforma de publicación beneficia a los estudiantes, las revistas selectivas ofrecen beneficios acumulativos superiores tanto para los estudiantes—por la retroalimentación recibida y el prestigio de publicar en una revista de calidad—como para la institución y la comunidad académica en general. Este hallazgo respalda la intención de incorporar en la presente propuesta un proceso de revisión por pares realizado por alumnos de grados superiores.

Finalmente, la experiencia de la Washington University Journal of Undergraduate Research (WUJUR) ilustra cómo una revista estudiantil puede operar con un proceso estructurado y formal de revisión por pares, colaborar con departamentos universitarios y otras revistas, y ofrecer a los estudiantes una vía transparente y gratuita para la publicación de sus trabajos \citep{wujur2024mission}. Este modelo resulta pertinente para la presente propuesta, pues comparte el objetivo de visibilizar la producción estudiantil y fomentar la curiosidad y la excelencia académica.

En síntesis, la evidencia revisada sustenta que la implementación de una revista científica interna en la carrera de Ingeniería en Biotecnología no solo contribuiría a la formación de los estudiantes en competencias de comunicación científica y pensamiento crítico, sino que también fortalecería su identidad académica y motivación, al tiempo que generaría un acervo documental de valor para la comunidad educativa.